% %%%%%%%%%%%%%%%%%%%%%%%%%%%%%%%%%%%%%%%%%%%%%%%%%%%%%%%%%%%%%%%%%%%%%%%%%%%%%%%%%
%
%	cap#.tex 
%  -----------
%
%	Neste arquivo você escreve o capítulo. Deve sempre iniciar com o ambiente
%	\chapter{Nome do Capitulo}
%
% %%%%%%%%%%%%%%%%%%%%%%%%%%%%%%%%%%%%%%%%%%%%%%%%%%%%%%%%%%%%%%%%%%%%%%%%%%%%%%%%%
%
\chapter{Análises considerando o cenário de túneis gêmeos}

Tendo em vista os assuntos abordados até aqui, este capítulo introduz o estudo de túneis gêmeos sem galeria transversal, e relaciona os mesmos com o túneis isolados já apresentados anteriormente. Dessa maneira, é possível analisar a influência da proximidade entre os túneis no comportamento do maciço e revestimento. As análises se concentram na convergência da seção do túnel ao longo do seu comprimento, nas pressões no revestimento e também na deformação ao longo da seção, que será dada por $\rho(\theta)$.

Os resultados referentes à convergência da seção do túnel $U=-u(r=R_t)/R_t$ e à pressão no revestimento $p=-\sigma_{rr}/c$, assim como no capítulo anterior, são apresentados no final da construção e no longo prazo, ao longo do comprimento do túnel, agora expresso por $z/R_t$. Para representar a convergência no equilíbrio, é dado o ponto $U_{eq}$ na coordenada $z=-140R_t$, ou seja, afastado da influência da frente de escavação. Da mesma forma, esse ponto também caracteriza as pressões no equilíbrio $p_{eq}$.

Diferentes pontos na periferia dos túneis foram analisados para a obtenção dos resultados, como demonstra a \autoref{desenho_pontos}, para ser possível analisar a influência do segundo túnel não apenas em um ponto específico da seção. Para os perfis de convergência e para as pressões no revestimento, neste capítulo serão apresentados somente os resultados referentes ao ponto $C(x=-d_1/2,y=R_t,z=-140R_t)$, relativo à coroa do túnel, e os demais estão presentes no Anexo A.

\begin{figure}[H]
	\begin{center}
		\includegraphics[scale=1.5]{desenho_pontos.pdf}
	\end{center}
	\caption{\label{desenho_pontos}Representação dos pontos de análise dos túneis gêmeos}
\end{figure}

\section{Influência da proximidade entre os túneis na convergência da seção}

A \autoref{TG_TI_UC_SR_inst} apresenta os resultados referentes ao final da construção, ou seja, na última escavação dos túneis, considerando o cenário de túneis não revestidos. Dessa maneira, é possível evidenciar o efeito da proximidade entre dois túneis longitudinais, sem a interferência dos efeitos da rigidez do revestimento, por exemplo. Nota-se um impacto expressivo na interação da proximidade entre os túneis longitudinais em comparação com um túnel único, representando um aumento de até 27,1\% na convergência, analisando o ponto C localizado na coroa.

\begin{figure}[H]
	\begin{center}
		\includegraphics[scale=0.8]{TG_TI_UC_SR_inst.pdf}
	\end{center}
	\caption{\label{TG_TI_UC_SR_inst}Perfis de convergência - túnel isolado (TI) e túneis gêmeos (TG) sem galeria (SG) e sem revestimento (SR) no final da construção}
\end{figure}

O aumento do parâmetro friccional promove uma redução significativa na convergência dos túneis gêmeos, uma vez que a discrepância observada, da ordem de 20,1\%, supera consideravelmente a variação identificada no túnel isolado, que foi de 11,6\%.

A \autoref{TG_TI_UC_SR_longo} apresenta perfis de convergência análogos aos do caso anterior, porém em estágio de longo prazo. Observa-se que, assim como no cenário de túnel isolado, o ângulo de atrito não exerce influência significativa sobre a convergência dos túneis gêmeos, não se verificando variações expressivas em função desse parâmetro. Isso sugere que, após a plastificação, a deformação em longo prazo é governada pela parcela viscoplástica, enquanto que o ângulo de atrito atua principalmente na fase plástica inicial.

\begin{figure}[H]
	\begin{center}
		\includegraphics[scale=0.8]{TG_TI_UC_SR_longo.pdf}
	\end{center}
	\caption{\label{TG_TI_UC_SR_longo}Perfis de convergência - túnel isolado (TI) e túneis gêmeos (TG) sem galeria (SG) e sem revestimento (SR) no longo prazo}
\end{figure}

Após verificar a influência da proximidade sem a interferência de um revestimento, é apresentado na \autoref{TG_TI_UC_CRE_inst} os perfis de convergência dos mesmos túneis (isolado e gêmeos) ao final da construção considerando a instalação de um revestimento elástico de concreto. A influência da proximidade, neste caso, é expressa por um aumento na convergência da seção em até 11,1\%, para o caso de $\phi=0^\circ$, sendo possível perceber que a utilização do revestimento elástico acarreta em uma limitação das convergências.

Uma redução de 34,8\% é observada na convergência do túnel com $\phi=0^\circ$ quando utilizado o revestimento elástico em comparação ao túnel não revestido. 


\begin{figure}[H]
	\begin{center}
		\includegraphics[scale=0.8]{TG_TI_UC_CRE_inst.pdf}
	\end{center}
	\caption{\label{TG_TI_UC_CRE_inst}Perfis de convergência - túnel isolado (TI) e túneis gêmeos (TG) sem galeria (SG), com revestimento elástico (CRE) no final da construção}
\end{figure}

Quando analisadas as curvas de convergência a longo prazo dos túneis revestidos com concreto elástico, percebe-se uma pequena influência dos diferentes ângulos de atrito para o maciço, influência essa não observada nos túneis não revestidos. Ainda, a utilização do revestimento limita a continuação das deformações, sendo um aumento de apenas 7,0\% em relação ao final da construção, para o caso de $\phi=10^\circ$.

\begin{figure}[H]
	\begin{center}
		\includegraphics[scale=0.8]{TG_TI_UC_CRE_longo.pdf}
	\end{center}
	\caption{\label{TG_TI_UC_CRE_longo}Perfis de convergência - túnel isolado (TI) e túneis gêmeos (TG) sem galeria (SG), com revestimento elástico (CRE) no longo prazo}
\end{figure}

\begin{figure}[H]
	\begin{center}
		\includegraphics[scale=0.8]{TG_TI_UC_CRVE_inst.pdf}
	\end{center}
	\caption{\label{TG_TI_UC_CRVE_inst}Perfis de convergência - túnel isolado (TI) e túneis gêmeos (TG) sem galeria (SG), com revestimento viscoelástico (CRVE) no final da construção}
\end{figure}

\begin{figure}[H]
	\begin{center}
		\includegraphics[scale=0.8]{TG_TI_UC_CRVE_longo.pdf}
	\end{center}
	\caption{\label{TG_TI_UC_CRVE_longo}Perfis de convergência - túnel isolado (TI) e túneis gêmeos (TG) sem galeria (SG), com revestimento viscoelástico (CRVE) no longo prazo}
\end{figure}

\section{Influência da proximidade entre os túneis nas pressões no revestimento}

\begin{figure}[H]
	\begin{center}
		\includegraphics[scale=0.8]{TG_TI_PC_CRE_inst.pdf}
	\end{center}
	\caption{\label{TG_TI_PC_CRE_inst}Pressões no revestimento - túnel isolado (TI) e túneis gêmeos (TG) sem galeria (SG), com revestimento elástico (CRE) no final da construção}
\end{figure}

\begin{figure}[H]
	\begin{center}
		\includegraphics[scale=0.8]{TG_TI_PC_CRE_longo.pdf}
	\end{center}
	\caption{\label{TG_TI_PC_CRE_longo}Pressões no revestimento - túnel isolado (TI) e túneis gêmeos (TG) sem galeria (SG), com revestimento elástico (CRE) no longo prazo}
\end{figure}

\begin{figure}[H]
	\begin{center}
		\includegraphics[scale=0.8]{TG_TI_PC_CRVE_inst.pdf}
	\end{center}
	\caption{\label{TG_TI_PC_CRVE_inst}Pressões no revestimento - túnel isolado (TI) e túneis gêmeos (TG) sem galeria (SG), com revestimento viscoelástico (CRVE) no final da construção}
\end{figure}

\begin{figure}[H]
	\begin{center}
		\includegraphics[scale=0.8]{TG_TI_PC_CRVE_longo.pdf}
	\end{center}
	\caption{\label{TG_TI_PC_CRVE_longo}Pressões no revestimento - túnel isolado (TI) e túneis gêmeos (TG) sem galeria (SG), com revestimento viscoelástico (CRVE) no longo prazo}
\end{figure}

\section{Anisotropia em relação a theta}

\begin{figure}[H]
	\begin{center}
		\includegraphics[scale=0.8]{TG_SG_rho_CRE.pdf}
	\end{center}
	\caption{\label{TG_SG_rho_CRE}Deformada da seção - túneis gêmeos sem galeria com revestimento elástico (CRE)}
\end{figure}

\begin{figure}[H]
	\begin{center}
		\includegraphics[scale=0.8]{TG_SG_rho_CRVE.pdf}
	\end{center}
	\caption{\label{TG_SG_rho_CRVE}Deformada da seção - túneis gêmeos sem galeria com revestimento viscoelástico (CRVE)}
\end{figure}

